\chapter{\ploren{Przegląd literatury i podstawy teoretyczne}{Literature review and theoretical background}}

Tytuł rozdziału musi dobrze odzwierciedlać jego treść, ale jednocześnie nie powinien być zbyt długi. Dobrze byłoby, gdyby jego długość nie przekraczała jednej linii. Na początku każdego rozdziału, przed podrozdziałem, powinny znaleźć się przynajmniej dwa, trzy zdania wprowadzenia, opisujące jego zawartość. Od nowej strony rozpoczynamy jedynie rozdziały główne; podrozdziały kontynuujemy na tej samej stronie. Oczywiście liczba rozdziałów pracy może być inna niż w niniejszym opracowaniu, inne będą także tytuły rozdziałów (z wyjątkiem dwóch pierwszych i ostatniego). W nagłówkach stron (z wyjątkiem strony tytułowej, spisu treści oraz pierwszych stron rozdziałów) umieszczamy tytuł rozdziału oraz numer strony.

W tym rozdziale autor m.in. opisuje koncepcje, modele, wzory matematyczne, które są kluczowe dla pracy, a także definiuje wykorzystywane w niej podstawowe pojęcia.

\section{Forma pracy dyplomowej}

Poniżej zademonstrowano formatowanie list numerowanych oraz wypunktowanych. Składniki list należy traktować jako elementy zdania, tzn. należy używać odpowiednich znaków interpunkcyjnych. Praca powinna spełniać następujące wymogi:

\begin{enumerate}
  \item format pracy – A4,
  \item kolejność stron w pracy:
\end{enumerate}

\begin{itemize}
  \item strona tytułowa,
  \item spis treści – generowany automatycznie,
  \item rozdział wprowadzający,
  \item rozdział zawierający cel i zakres pracy,
  \item rozdziały teoretyczne (opisowe) oraz praktyczne,
  \item rozdział podsumowujący pracę,
  \item wykaz bibliografii,
  \item ewentualne załączniki.
\end{itemize}

\begin{enumerate}
  \item Krój czcionki tekstu – Times New Roman; rozmiar – 12 pkt; rozmiar czcionki podpisów tabel, rysunków oraz listingów – 11 pkt.
  \item Rozmiar marginesów: lewego – 3,5 cm; prawego, górnego i dolnego – 2,5 cm.
  \item Formatowanie akapitów: interlinia 1,5 wiersza; justowanie obustronne; odległość między akapitami: 6 pt; wcięcie akapitowe: 0,75 cm (z wyjątkiem pierwszego akapitu w rozdziale lub podrozdziale – ten pozostawiamy bez wcięcia).
  \item Dalsze instrukcje dotyczące przygotowania pracy do obrony znajdują się na stronie https://cos.po.edu.pl.
\end{enumerate}

\section{Technika pisania pracy}

Podrozdziały nie powinny być zbyt krótkie, dlatego tutaj powinno znaleźć się więcej zdań.

\subsection{Numeracja rozdziałów}

Rozdziały numeruje się do trzeciego poziomu, tzn. 1., 1.1, 1.1.1, 1.2., 1.2.1., … itd. Rozdziały i podrozdziały powinny tworzyć logiczną, spójną hierarchię. Nie tworzymy pojedynczych podrozdziałów, np. podrozdział 4.1. ma sens tylko wtedy, gdy istnieje także podrozdział 4.2. oraz ewentualnie kolejne.

\subsection{Numeracja rysunków}

Wszystkie rysunki powinny mieć szczegółowe podpisy i powinny być numerowane dwoma cyframi rozdzielonymi kropką. Pierwsza cyfra oznacza numer głównego rozdziału, zaś druga – numer kolejny rysunku w rozdziale. Podpisy pod rysunkami powinny być wyśrodkowane. Po podpisie, podobnie jak po tytułach rozdziałów i podrozdziałów, nie stawiamy kropki. Odległość między akapitem a rysunkiem: 12 pt; między rysunkiem a jego podpisem: 6 pt; między podpisem rysunku a kolejnym akapitem: 12 pt. To samo dotyczy tabel, listingów oraz ich podpisów.

Rys. 3.1. Schemat uczenia nienadzorowanego sieci neuronowej w zadaniu kompresji obrazu [2]

Rysunek musi być czytelny i mieścić się w granicach tekstu; nie może wychodzić na margines. Należy zastosować układ formatowania „Równo z tekstem”. W przypadkach dużych, szerokich rysunków zastosować należy poziomą orientację strony. Rysunki należy wykonywać samodzielnie; w wyjątkowych przypadkach (np. bardzo złożony schemat) dopuszczalne jest wykorzystywanie „skanów” z książek lub grafik ze źródeł internetowych. Rysunki przedstawiające „zrzuty” okien programów nie powinny być zbyt duże (estetyka) ani zbyt małe (czytelność), ich rozmiar powinien być jednakowy w całej pracy. Jeżeli rysunek pochodzi z innego źródła, odnośnik do tego źródła należy umieścić po podpisie w nawiasach kwadratowych, analogicznie jak w przypadku odnośników znajdujących się w tekście pracy.

Rys. 3.2. Wpływ wartości współczynnika przenikania ciepła na przebieg reakcji(opracowanie własne): a) Zmiany temperatury b) Zmiany zawartości produktu reakcji

Do każdego zamieszczonego w pracy rysunku powinno być przynajmniej jedno odwołanie z jej treści. Odwołujemy się wyłącznie poprzez numer rysunku, nigdy z użyciem słów „powyżej”, „poniżej”. Na rys. 3.1. przedstawiono schemat ilustrujący proces uczenia nienadzorowanego sieci neuronowej, wykorzystywanej w zadaniu stratnej kompresji obrazu. Rysunek 3.2.a) przedstawia wykres ilustrujący wpływ wartości współczynnika przenikania ciepła U na… Z kolei na rys. 3.2.b) zamieszczono wykres pokazujący…

\subsection{Wzory}

Jeżeli w pracy występują wzory, do ich wpisywania należy użyć edytora równań. Niedopuszczalne jest wstawianie wzorów w postaci skanów bądź rysunków. Wzory wyrównujemy do środka (tabulator) i numerujemy w analogiczny sposób jak rysunki. Numer wzoru umieszczamy w nawiasie, z prawej strony przy marginesie. Rozmiar znaków we wzorach powinien odpowiadać rozmiarowi w tekście pracy. Wzory należy traktować jako element zdania, tzn. należy używać po nich znaków interpunkcyjnych. Pod wzorem wyjaśniamy wszystkie występujące w nim oznaczenia, chyba, że zostały wprowadzone już wcześniej. Przyrosty współczynników wagowych $\Delta w(i)$ wyznaczane są w i-tym kroku uczenia sztucznego neuronu zgodnie z następującą zależnością:

(3.1)gdzie:

– wektor współczynników wagowych w i-tym kroku uczenia,

– gradient funkcji błędu sieci,

– współczynnik prędkości uczenia (0-1) w i-tym kroku uczenia,

$\alpha$ – stała momentum (0-1).

Odwołania do wzorów z treści pracy oraz wtrącenia obcojęzyczne powinny wyglądać jak w kolejnym zdaniu. Zależność (3.1) opisuje modyfikację współczynników wagowych neuronu przy użyciu GDM (ang. gradient descent with momentum), czyli algorytmu największego spadku uwzględniającego tzw. człon momentum (pędu). Każdy akronim czyli skrótowiec musi być wyjaśniony (rozwinięty) przy jego pierwszym użyciu. W dalszej części pracy taki akronim należy stosować już bez jego rozwijania.

Symbole wielkości fizycznych piszemy kursywą, zaś ich wartości liczbowe oraz symbole jednostek – czcionką prostą. Pamiętamy przy tym, że wartość liczbowa i symbol jednostki wielkości fizycznej muszą być oddzielone pojedynczą spacją, np. Usk = 230 V. Symbole oraz argumenty funkcji, takie jak np. f(x), (), piszemy czcionką pochyłą, zaś nawiasy – czcionką prostą. Symbole stałych matematycznych, takich jak e oraz $\pi$, a także symbole nazw funkcji matematycznych, takie jak sin x oraz log x, piszemy czcionką prostą.

\subsection{Tabele}

Tabele numerujemy tak jak wzory, jednak podpis umieszczamy nad tabelą (bez kropki na końcu). Należy pamiętać o tym, żeby w całej pracy tabele miały podobny wygląd (rodzaj czcionki, ewentualne pogrubienia w nagłówku, itp.).

\textbf{Tab. 3.1. Charakterystyka poszczególnych warstw standardu MPEG-1}

Rozdziału lub podrozdziału nigdy nie rozpoczynamy ani kończymy rysunkiem, tabelą, wzorem. Powinny zawsze znaleźć się tu przynajmniej 2-3 zdania wprowadzenia lub podsumowania. W omawianym przypadku będzie to komentarz do tab. 3.1.

\subsection{Spacje i znaki interpunkcyjne}

W pracy nie należy używać wielokrotnych spacji, ani wielokrotnych znaków nowego akapitu. Znaki interpunkcyjne, takie jak przecinek (,), kropka (.), dwukropek (:), średnik (;), znak zapytania (?), wykrzyknik (!), zamknięcie dowolnego nawiasu (]\})>), zamknięcie cudzysłowu (” lub ’) nie mogą być poprzedzone spacją. Bezpośrednio po wymienionych znakach może wystąpić wyłącznie spacja, znak nowego akapitu lub inny znak interpunkcyjny. Po znakach otwierających dowolnego nawiasu ([\{(<) lub cudzysłowu („ lub ') nigdy nie należy używać spacji. Spację używamy przed tymi znakami. Nie należy pozostawiać spacji na końcu akapitu – przed znakiem nowego akapitu. Tytuły rozdziałów i podrozdziałów, a także rysunków, tabel, listingów pozostawiamy bez kropki na końcu.

Pamiętamy, że w roli myślnika nie używamy łącznika ani znaku „minus” (-) lecz tzw. półpauzy (–). Znak ten można uzyskać, wciskając jednocześnie klawisz <Ctrl> i minus z klawiatury numerycznej. Przykłady użycia łącznika oraz półpauzy zademonstrowano np. w rozdziale 3.2.3, w opisie wielkości pod wzorem (3.1).

Apostrof stawiamy wyłącznie przy tych wyrazach, których końcówka jest niema. Jeśli w wymowie nie opuszcza się żadnej głoski, jego użycie jest niepoprawne, np. John’a, Oliver’a. W obu przypadkach powinno stosować się tutaj zapis fonetyczny, tj. Johna, Olivera. Natomiast poprawnie będzie np.: skorzystać z pendrive’a, zastosować transformację Laplace’a, omówić prawo Joule’a–Lenza.

\subsection{Kod programu lub skryptu}

Krótki fragment kodu programu, procedury lub skryptu piszemy używając czcionki maszynowej (Courier lub Consolas) o rozmiarze 11 pt:

\begin{lstlisting}
wynik = add(10,20);/* przyklad 1 */
wynik = add(x1*20+4,x1/x2);/* przyklad 2 */
\end{lstlisting}

Inną możliwością prezentacji dłuższego kodu, np. całej procedury, skryptu, programu lub wyników ich działania, jest użycie tabeli (ramki) w postaci zaprezentowanej na listingu 3.1. Zdecydowanie preferowane jest zamieszczanie kodu w postaci tekstowej. W niektórych przypadkach (np. gdy zależy nam na pokazaniu całego okna programu lub gdy nie jest możliwe łatwe skopiowanie kodu w postaci tekstowej) można użyć listingu w formie graficznej, np. „zrzutu” okna programu. Jednak podpisujemy go jako listing, a nie jako „rysunek” – liczy się forma logiczna, a nie techniczna. Ta sama uwaga dotyczy np. tabel wstawianych w formie graficznej – są to tabele, a nie „rysunki”.

\textbf{List. 3.1. Prosty program w języku C}

Należy unikać pozostawiania pojedynczych liter (np. spójników „w”, „i”, „o”, „z”) na końcu wierszy, zarówno w standardowym tekście rozdziału, jak i w nagłówkach, podpisach rysunków, itp. W tym celu można wykorzystać tzw. twardą spację (ang. hard space) – znak wyglądający tak samo jak zwykła spacja, lecz przez edytor traktowany jak litera lub cyfra, czyli część wyrazu. W programie Microsoft Word znak twardej spacji wstawia się przy pomocy kombinacji klawiszy <Ctrl><Shift><spacja>. Istnieje także twardy myślnik (<Ctrl><Shift><->), który ma działanie podobne do twardej spacji – sprawia, że myślnik staje się częścią wyrazu.

\subsection{Wyróżnienia w tekście}

Wyróżnianie niektórych elementów tekstu poprawia jego czytelność. Ważne jest, żeby w całej pracy jednakowe elementy były oznaczane tak samo. Poniżej przedstawiono przykłady wyróżnień:

\begin{itemize}
  \item nazwy programów (notepad.exe), aplikacji (Android Studio), systemów operacyjnych (Microsoft Windows), itp., również piszemy kursywą,
  \item nowe, mało znane albo bardzo istotne pojęcia lub sformułowania można wyróżnić pogrubieniem (jednak nie należy nadużywać takich wyróżnień),
  \item elementy języka programowania – zarówno występujące w tekście pracy, jak i w formie listingu, wyróżniać należy czcionką maszynową (Courier lub Consolas).
\end{itemize}

Unikamy wielokrotnych wyróżnień tekstu. Jedna z podstawowych zasad typografii mówi, że nie należy wyróżniać fragmentu tekstu na więcej niż jeden sposób, jeśli nie jest to konieczne. W szczególności zasada ta dotyczy jednoczesnego stosowania pogrubienia i kursywy. Inne przykłady sytuacji, których należy unikać, to użycie pochylenia i cudzysłowów w cytatach oraz podkreślenie i kolorowanie linków. Tekst bez podwójnych wyróżnień jest bardziej estetyczny i nie rozprasza czytelnika.

Nowy rozdział rozpoczynamy od nowej strony – tylko w takim przypadku wewnątrz pracy może pojawić się „pusty obszar” jak poniżej.
